% Drawing 74/04-5 -- folder 74, batch 04, archive page 66 (PDF page 67).
% Transcription: un hexagone à sommet central $b$ ; sommets étiquetés, en partie surchargés : en haut à gauche $b' = A_4$, en haut à droite $\cdot = A_3$, à gauche $a'$ (précédé de « 2: »), à droite $\cdot = A_2$, en bas à gauche $c = A_0$, en bas à droite $a = A_1$ ; les arêtes $a'b'$, $A_0A_1$ et $A_1A_2$ sont repassées en gras. À côté, deux petits schémas croisant des lignes entre $\xi_1, \xi_2, \xi_3$ et $\xi_1^\wedge, \xi_2^\wedge, \xi_3^\wedge$ (lecture douteuse).
% Kind: hexagon with central vertex; two small crossing schemas. Method 2 (schematic redraw from the facsimile). Status: partial.
% Notes: Bbox widened to the right and down (from [0.12, 0.06, 0.75, 0.2]) so that the second small schema is whole; all coordinates are in units of the new crop, the figure drawn at 2.4 mm per unit so that the \small labels keep their places. Hexagon with the sides a'b' (left upper side), A0A1 (bottom) and A1A2 (right lower side) bold; inside it, as on the page, not two crossing diagonals but a stroke from the top right vertex to the central point b and on to the bottom left vertex, a stroke from b to the bottom right vertex, and a stroke from the top left vertex that stops short of b; the points b and A2 are blots. Labels: b' = A4 top left, . = A3 top right, 2: a' left, . = A2 right, c = A0 bottom left, a = A1 bottom right, b beside the centre. Beside, two small schemas: each a row xi1 xi2 xi3 above and a row xi1^ xi2^ xi3^ below (hats read as in the transcription, doubtful), joined by crossing strokes: in the first, three long strokes and a short one; in the second, a small circle at the upper left from which a long stroke runs to the lower right, a stroke from the lower left up to the foot of xi2 and down again to the right, a long stroke from the lower left to the upper right, a shorter one parallel to the last, a short stroke under the circle, and a bracket under xi2^. Redrawn 2026-11-09 (issue 128); overlay check 2026-11-13 (issue 147). Second pass 2026-11-14 (issue 147, Claude claude-opus-5-5): every stroke re-read on 22 to 60 px/unit gridded crops of the 400 dpi colour scan and redrawn there (the hexagon's inner strokes were drawn as two diagonals crossing at a computed centre, and the schemas' strokes were approximate). Rasterised by the session's previewer (not compiled: no TeX in the session) and laid over the crop: every stroke lies on its ink to within about 1 unit. Still differs: the labels of the two small schemas are a doubtful reading (as in the transcription): the lower row's marks may be hats, primes or indices, and the bracket under xi2^ is not explained. Nothing is struck. Not checked by a reader.
% Coordinates: 1 unit = 2.4 mm = 1/100 of the crop width (bbox in page fractions: [0.12, 0.06, 0.88, 0.23]).
% Redrawn by Claude (claude-opus-5-5), issue #73, 2026-11-09; not checked by a human; not compiled (no TeX in the session).
% Checked by Claude (claude-opus-5-5) 2026-11-13 (issue #147): overlaid on the 300 dpi colour facsimile; not compiled (no TeX in the session). Status: partial.
% Second pass by Claude (claude-opus-5-5) 2026-11-14 (issue #147): redrawn stroke by stroke from the 400 dpi colour facsimile, overlaid on a grid; not compiled (no TeX in the session). Status: partial.
% Rights: reproduces part of the work -- behind RENDER_TRANSCRIPTS, never published by itself.
\documentclass[tikz,border=3pt]{standalone}
\usepackage{amsmath,amssymb}
\usetikzlibrary{patterns,arrows.meta}
\tikzset{pencil/.style={black!45}, struck/.style={black!35}, mark/.style={circle,draw,inner sep=1.3pt}, pt/.style={circle,fill,inner sep=0.9pt}, opt/.style={circle,draw,fill=white,inner sep=1.1pt}}
\begin{document}
\begin{tikzpicture}[x=2.4mm,y=2.4mm,line cap=round,line join=round,font=\scriptsize]
\draw[very thick] (11.03,-15.34) -- (16,-9.95);
\draw (16,-9.95) -- (24.62,-10.11) -- (28.76,-16.16);
\draw[very thick] (28.76,-16.16) -- (24.62,-21.8) -- (15.92,-21.8);
\draw (15.92,-21.8) -- (11.03,-15.34);
\draw (24.62,-10.11) -- (19.73,-16.99) -- (15.92,-21.8);
\draw (19.73,-16.99) -- (24.62,-21.8);
\draw (16.2,-10.4) -- (18.9,-16.41);
\fill (19.73,-16.99) circle[radius=0.3];
\fill (28.76,-16.16) circle[radius=0.3];
\node[inner sep=1pt] at (18.07,-8.12) {\small $b'=A_4$};
\node[inner sep=1pt] at (28.02,-9.45) {\small $\cdot=A_3$};
\node[inner sep=1pt] at (7.05,-15.34) {\small $2{:}\ a'$};
\node[inner sep=1pt] at (32.33,-16.25) {\small $\cdot=A_2$};
\node[inner sep=1pt] at (19.07,-23.71) {\small $c=A_0$};
\node[inner sep=1pt] at (26.36,-23.54) {\small $a=A_1$};
\node[inner sep=1pt] at (21.39,-16) {\small $b$};
% first small schema
\draw (57.1,-18.2) -- (61.5,-19.6);
\draw plot[smooth] coordinates {(55.9,-20) (58,-19.3) (62,-18.4)};
\draw (58.9,-19.75) -- (60.4,-19) -- (62,-18.4);
\draw (57.25,-18.5) -- (58.5,-19.9);
\node[inner sep=1pt] at (56.2,-16.8) {\small $\xi_1$};
\node[inner sep=1pt] at (59.3,-16.7) {\small $\xi_2$};
\node[inner sep=1pt] at (63.3,-17) {\small $\xi_3$};
\node[inner sep=1pt] at (56.7,-22) {\small $\hat\xi_1$};
\node[inner sep=1pt] at (60.3,-22) {\small $\hat\xi_2$};
\node[inner sep=1pt] at (63.7,-22) {\small $\hat\xi_3$};
% second small schema
\draw (80.5,-21.1) circle[radius=0.35];
\draw (80.8,-21.25) -- (91.25,-25.1);
\draw (79.4,-23.75) -- (84.75,-21.1) -- (90,-24.5);
\draw (79.4,-23.75) -- (89.75,-21.4);
\draw (84,-24.5) -- (90.5,-21.4);
\draw (81.75,-22.5) -- (83.1,-23.9);
\draw (84.1,-27.4) -- (84.1,-29.1) -- (88.9,-29.1) -- (88.8,-27.4);
\node[inner sep=1pt] at (77.8,-19) {\small $\xi_1$};
\node[inner sep=1pt] at (85.5,-18.75) {\small $\xi_2$};
\node[inner sep=1pt] at (92.75,-18.75) {\small $\xi_3$};
\node[inner sep=1pt] at (78.4,-27.6) {\small $\hat\xi_1$};
\node[inner sep=1pt] at (86.5,-27.6) {\small $\hat\xi_2$};
\node[inner sep=1pt] at (94,-27) {\small $\hat\xi_3$};
\end{tikzpicture}
\end{document}
